\documentclass[10pt]{article}
\usepackage[margin=1in]{geometry}
\usepackage{booktabs}
\usepackage{multirow}
\usepackage{amsmath}
\usepackage{array}
\usepackage{placeins}
\usepackage{graphicx}
\usepackage{hyperref}

\title{Where Does PL Actually Beat Deterministic?\\A Query-Level Investigation of the EE-D-Interval Table}
\author{}
\date{}

\begin{document}
\maketitle

\noindent\textbf{This is a working diagnostic document, not part of the paper.} The
paper's EE-D-interval table (Table~\ref{tab:baseline-repro} below, reproducing
\texttt{tab:baseline} in \texttt{fair\_rag\_diversity\_paper.tex}) bins PL's utility
gain over deterministic by EE-D disparity interval, pooled across all four
$\alpha$, for the two tasks it covers (LaMP-1 and LaMP-4, the only tasks with a
published source-paper baseline to compare against). Some of those bins are
\emph{positive} --- PL beats deterministic on average. This document asks: what is
actually different about the queries in those bins? Specifically:
\begin{enumerate}
  \item Is it just a handful of queries and sampling noise?
  \item Did deterministic's top-5 already contain an unusually high fraction of
    genuinely useful documents?
  \item Did deterministic fail to retrieve anything useful, leaving PL nowhere to
    go but up?
  \item Does MMR show the same pattern on the identical queries, or is this
    specific to PL?
\end{enumerate}
\textbf{No new experiments were run.} Everything here is computed from already-completed
runs under \texttt{experiment\_runs/}, using the same with-gold EU$_{\text{norm}}$
normalization (\texttt{analysis.with\_gold\_normalization}) as every other table in
this project. ``Useful'' document means a profile item with
\texttt{relevance\_label}${}=1$ in this project's utility-delta relevance mapping
(\texttt{data/lamp\_utility\_labels\_*/*\_relevance\_mapping.tsv}, built by
\texttt{utility\_labels/analyze\_delta.py}) --- i.e.\ a document that
empirically changed that generator's output when swapped in or out, \emph{not}
topical similarity. \textbf{Notation:} $^*p{<}0.05$, Welch's $t$-test vs.\ the
cell's full pooled deterministic set, same convention as the main paper.

\section{The Table This Investigates}
\label{sec:source-table}

\begin{table}[h]
\centering
\scriptsize
\begin{tabular}{@{}lllccccc@{}}
\toprule
Task & Gen. & Ranker & [.0,.2) & [.2,.4) & [.4,.6) & [.6,.8) & [.8,1.0] \\
\midrule
LaMP-1 & Base & BM25 & $-$0.23* (n{=}341) & $-$0.09* (n{=}131) & $-$0.05 (n{=}42) & $-$0.05 (n{=}106) & $-$0.02 (n{=}296) \\
LaMP-1 & Base & Contr. & $-$0.15* (n{=}363) & $+$0.01 (n{=}91) & $-$0.07 (n{=}63) & $+$0.05 (n{=}132) & $+$0.01 (n{=}255) \\
LaMP-1 & Small & BM25 & $-$0.09 (n{=}68) & $-$0.14 (n{=}1) & $+$0.01 (n{=}7) & $+$0.03 (n{=}17) & $-$0.01 (n{=}43) \\
LaMP-1 & Small & Contr. & $-$0.05 (n{=}60) & $-$0.15* (n{=}9) & $-$0.10 (n{=}11) & $+$0.14 (n{=}17) & $-$0.02 (n{=}39) \\
LaMP-4 & Base & BM25 & $-$0.08* (n{=}1040) & $-$0.02 (n{=}379) & $+$0.02 (n{=}347) & $+$0.00 (n{=}484) & $-$0.00 (n{=}1054) \\
LaMP-4 & Base & Contr. & $-$0.13* (n{=}1080) & $-$0.04* (n{=}336) & $-$0.00 (n{=}402) & $-$0.01 (n{=}463) & $+$0.00 (n{=}1023) \\
LaMP-4 & Small & BM25 & $-$0.11* (n{=}142) & $-$0.02 (n{=}47) & $-$0.01 (n{=}34) & $-$0.05 (n{=}71) & $-$0.00 (n{=}838) \\
LaMP-4 & Small & Contr. & $-$0.10* (n{=}1110) & $-$0.04* (n{=}335) & $-$0.01 (n{=}400) & $+$0.02 (n{=}450) & $+$0.00 (n{=}1037) \\
\bottomrule
\end{tabular}
\caption{$\Delta$EU$_{\text{norm}}$ vs.\ deterministic by EE-D interval, pooled
across all four PL $\alpha$ (reproducing the paper's Table~\ref{tab:baseline-repro}
body, with $n$ --- number of (qid,$\alpha$) observations in that bin --- added).
Matches \texttt{report/tables/pl\_improve\_bin\_summary.csv}.}
\label{tab:baseline-repro}
\end{table}

Of these 40 (cell, bin) combinations, \textbf{11 are positive and 29 are negative}.
\textbf{Zero of the 11 positive bins are statistically significant} (all
$p\geq0.17$); all 10 significant bins in the table are negative. So in the strict
sense, no cell of this table contains a defensible instance of ``PL significantly
beats deterministic'' --- every positive number here is, by the table's own
significance test, indistinguishable from noise. That is itself the headline
answer to question 1, and the rest of this document looks at whether the noise is
at least \emph{structured} --- landing preferentially on a particular kind of
query rather than being uniform.

\FloatBarrier
\section{What's Different About the Queries in the Positive Bins}
\label{sec:detail}

Table~\ref{tab:detail} gives the full detail for all 11 positive bins, sorted by
effect size. Two columns matter most: \texttt{cell\_relevant} is the average
fraction of useful documents in deterministic's top-5, averaged over \emph{every}
query in that cell (all 5 bins together); \texttt{bin\_relevant} is that same
average restricted to \emph{just} the queries that landed in this specific bin.
Same idea for \texttt{cell\_det\_EU}/\texttt{bin\_det\_EU}, but for deterministic's
own EU$_{\text{norm}}$ instead of retrieval quality.

\begin{table}[h]
\centering
\tiny
\begin{tabular}{@{}lllccrrrrrrrr@{}}
\toprule
Task & Gen. & Ranker & Bin & $n_{\text{obs}}$/$n_{\text{qid}}$ & $\Delta_{\text{PL}}$ & $p$ & cell\_relevant & bin\_relevant & bin\_PL\_relevant & cell\_det\_EU & bin\_det\_EU & $\Delta_{\text{MMR}}^{\text{same qids}}$ \\
\midrule
LaMP-1 & Small & Contr. & [.6,.8) & 17/17 & $+$0.144 & 0.19 & 0.235 & 0.341 & 0.326 & 0.147 & 0.294 & $+$0.012 \\
LaMP-1 & Base & Contr. & [.6,.8) & 132/132 & $+$0.054 & 0.25 & 0.372 & 0.382 & 0.375 & 0.633 & 0.674 & $-$0.017 \\
LaMP-1 & Small & BM25 & [.6,.8) & 17/16 & $+$0.029 & 0.79 & 0.235 & 0.362 & 0.330 & 0.176 & 0.188 & $-$0.062 \\
LaMP-4 & Base & BM25 & [.4,.6) & 347/281 & $+$0.021 & 0.17 & 0.356 & 0.389 & 0.395 & 0.290 & 0.320 & $+$0.007 \\
LaMP-4 & Small & Contr. & [.6,.8) & 450/409 & $+$0.019 & 0.19 & 0.361 & 0.355 & 0.377 & 0.326 & 0.320 & $+$0.002 \\
LaMP-1 & Small & BM25 & [.4,.6) & 7/7 & $+$0.014 & 0.93 & 0.235 & 0.257 & 0.240 & 0.176 & 0.286 & $-$0.057 \\
\addlinespace
LaMP-1 & Base & Contr. & [.8,1.0] & 255/225 & $+$0.009 & 0.84 & 0.372 & 0.380 & 0.386 & 0.633 & 0.631 & $-$0.004 \\
LaMP-1 & Base & Contr. & [.2,.4) & 91/77 & $+$0.007 & 0.87 & 0.372 & 0.418 & 0.324 & 0.633 & 0.636 & $-$0.010 \\
LaMP-4 & Base & BM25 & [.6,.8) & 484/443 & $+$0.005 & 0.77 & 0.356 & 0.356 & 0.371 & 0.290 & 0.281 & $+$0.004 \\
LaMP-4 & Small & Contr. & [.8,1.0] & 1037/827 & $+$0.002 & 0.86 & 0.361 & 0.360 & 0.378 & 0.326 & 0.325 & $-$0.003 \\
LaMP-4 & Base & Contr. & [.8,1.0] & 1023/820 & $+$0.002 & 0.89 & 0.434 & 0.434 & 0.441 & 0.339 & 0.338 & $-$0.004 \\
\bottomrule
\end{tabular}
\caption{All 11 positive bins, full detail. $n_{\text{obs}}$ is (qid,$\alpha$)
observations, $n_{\text{qid}}$ is unique queries (they differ when one query lands
in the same bin at two different $\alpha$'s). Split by a blank row into the 6 bins
with a meaningful effect size ($\Delta_{\text{PL}}\geq0.01$) and the 5 that are
themselves near-zero. Full data in \texttt{pl\_improve\_bin\_detail}.csv.}
\label{tab:detail}
\end{table}

\begin{figure}[h]
\centering
\includegraphics[width=\textwidth]{fig_pl_improvement_bins.png}
\caption{Left: PL's own delta vs.\ MMR's delta on the identical queries. Middle:
fraction of useful documents in deterministic's top-5, this bin vs.\ the cell-wide
average. Right: deterministic's own EU$_{\text{norm}}$, this bin vs.\ the cell-wide
average. Rows ordered by PL's effect size (same order as Table~\ref{tab:detail}).}
\label{fig:bins}
\end{figure}

\FloatBarrier
\subsection*{Question 1: is it just a few queries and noise?}
Partially, and it's worth separating two different things sample size affects.
Some of the 11 bins genuinely are tiny --- 17, 17, and 7 queries in the three
smallest --- and at those sizes a $+0.03$--$+0.14$ swing is well within normal
sampling variation for a per-query metric this noisy. But the effect-size split in
Table~\ref{tab:detail} does \emph{not} track sample size: the largest bin
(LaMP-4/Base/Contr., $n{=}1023$ unique queries) has one of the smallest effects
($+0.002$), while a mid-sized bin (LaMP-1/Base/Contr., $n{=}132$) has the
second-largest ($+0.054$). And critically, \textbf{even the large-$n$ bins never
reach significance} --- $n{=}1023$ and $p{=}0.89$ is not ``needs more data,'' it's
``the true effect at this size is close to zero.'' So the honest reading is: the 5
near-zero bins ($\Delta_{\text{PL}}{<}0.01$) are just noise around zero, full stop,
regardless of their $n$; the 6 larger bins are \emph{not} explained by small
samples alone (three of them have $n{\geq}132$), but none of them clears
significance either, so ``noise'' remains the statistically correct label even
there --- the question below is whether that noise is at least non-random in
\emph{where} it lands.

\subsection*{Question 2: was the \% of useful documents relatively high?}
Yes, cleanly, for the 6 bins with a real effect size. In \textbf{all 6}, deterministic's
own top-5 already contained a higher fraction of useful documents than that cell's
average query (e.g.\ LaMP-4/Base/BM25: 0.389 in-bin vs.\ 0.356 cell-wide; LaMP-1/Small/Contr.:
0.341 vs.\ 0.235 --- close to 50\% relatively higher). This pattern disappears for the
5 near-zero bins (only 2 of 5 show any elevation, and it's marginal in both).
Notably, PL's \emph{own} retrieved lists in the effect-size bins (\texttt{bin\_PL\_relevant})
are usually about the same as deterministic's, not higher --- PL is not
systematically finding more useful documents than deterministic already found for
these queries. The elevation is a property of \emph{which queries land in this EE-D
range}, not something PL's reranking adds on top. A plausible mechanism: queries
with more useful documents scattered through the candidate pool (rather than
concentrated at the very top) are the queries where repeated PL sampling disagrees
with itself the most --- which is exactly what raises EE-D --- and where swapping
one adequate document for another adequate one costs little to nothing, so PL's
variance nets out roughly neutral-to-positive instead of costly.

\subsection*{Question 3: did deterministic just not find anything useful in top-5?}
No --- if anything, close to the opposite. In 5 of the 6 effect-size bins,
deterministic's own EU$_{\text{norm}}$ in that bin is \emph{higher} than the
cell-wide average (LaMP-1/Small/Contr.: 0.294 vs.\ 0.147, roughly double;
LaMP-1/Base/Contr.: 0.674 vs.\ 0.633; LaMP-4/Base/BM25: 0.320 vs.\ 0.290). The one
exception, LaMP-4/Small/Contr.\ [.6,.8), is close to flat (0.320 vs.\ 0.326, a
trivial difference). So this hypothesis is not supported: PL's positive bins are
\emph{not} where deterministic was struggling and had nowhere to go but up. They
skew toward queries where deterministic was already doing relatively well, and PL
edges it further --- a genuinely counter-intuitive result relative to the
hypothesis, worth flagging explicitly rather than quietly not mentioning it.

\FloatBarrier
\section{Does MMR Show the Same Pattern on the Same Queries?}
\label{sec:mmr-compare}

The rightmost column of Table~\ref{tab:detail} and the left panel of
Figure~\ref{fig:bins} answer this directly: MMR's $\Delta$EU$_{\text{norm}}$ on the
\emph{identical} query sets (pooled across MMR's own 5 $\lambda$'s, same
with-gold normalization), compared to PL's on those same queries.

\textbf{MMR does not replicate PL's pattern.} Across the 11 bins, MMR is negative
in 7 and positive in only 4; averaged over all 11, PL's mean delta is $+0.028$
while MMR's is $-0.012$ --- opposite sign. In the three largest-effect bins
specifically, MMR is negative in two of three and, even where positive
(LaMP-1/Small/Contr.: $+0.012$), an order of magnitude smaller than PL's $+0.144$
on the same 17 queries. This means whatever is happening in these bins is not
simply ``these particular queries are easy to improve on regardless of method''
--- it is specific to PL's stochastic, resample-many-times mechanism. MMR's
deterministic greedy selection, applied to the same queries, mostly does not
capture it and sometimes loses ground where PL gains. This is consistent with
this project's other finding (Section 3 of \texttt{aux\_diagnostics.pdf}) that
PL's positive bins are well explained by ``which setting dominates the bin,''
while MMR's effects are compositionally distinct --- here we see the same
asymmetry from the opposite direction: PL's own positive bins do not transfer to
MMR even on identical queries.

\FloatBarrier
\section{Summary}
\label{sec:summary}

\begin{itemize}
  \item \textbf{None of the 11 positive EE-D-interval bins for PL are
    statistically significant} ($p\geq0.17$ throughout); every significant result
    in the table is negative. Read strictly, the table contains zero evidence
    that PL beats deterministic anywhere.
  \item Sample size explains the 5 smallest-effect bins (all $\Delta_{\text{PL}}{<}0.01$,
    indistinguishable from zero) but \emph{not} the 6 larger ones --- three of
    those have $n{\geq}132$ unique queries and still show a real point estimate,
    just not enough to clear significance at this project's scale.
  \item Those 6 larger-effect bins consistently over-represent queries where
    deterministic's own top-5 already had an above-average fraction of useful
    (utility-delta-labeled) documents, and where deterministic's own utility was
    already above the cell average --- the opposite of ``deterministic failed to
    find anything useful.'' PL is not adding more useful documents than
    deterministic found; it is landing, by chance across its samples, on
    comparably-useful alternatives for queries where several adequate options
    exist.
  \item MMR does not show this pattern on the same queries --- it is specific to
    PL's sampling mechanism, not a general property of these queries being
    ``easy.''
\end{itemize}

The overall picture: PL's apparent wins in the EE-D-interval table are real point
estimates but statistically indistinguishable from noise, and that noise is not
uniformly distributed --- it concentrates on relevance-rich queries where
deterministic was already relatively strong. That is a suggestive characterization
of \emph{where} PL's variance happens to land favorably, not evidence of a
systematic mechanism by which PL improves retrieval quality.

\end{document}
